% !PRINTED 404
\weberpage{404}
We can, however, further assume that $a_0 = 1$; for if in the product $a_0^{n-1} f(x)$ we set
\[
a_0 x = x_1,
\]
the equation becomes
\[
x_1^n + a_1 x_1^{n-1} + a_0 a_2 x_1^{n-2} + \cdots + a_0^{n-1} a_n = 0.
\]
We shall therefore now assume that in the equation
\weq{1}{x^n + a_1 x^{n-1} + a_2 x^{n-2} + \cdots + a_{n-1} x + a_n = 0}
the coefficients $a_1, a_2 \ldots a_n$ are integers. A \gesperrt{rational} root of (1) cannot be a fractional number; for if we assume that (1) is satisfied by the fraction
\[
x = \frac{p}{q},
\]
where $p$ and $q$ are integers without a common divisor, and $q$ is positive and greater than 1, then multiplication by $q^{n-1}$ gives
\[
\frac{p^n}{q} + a_1 p^{n-1} + a_2 q p^{n-2} + \cdots + a_n q^{n-1} = 0.
\]
Thus $p^n : q$ would have to be an integer, which is impossible.

Now if the equation (1) can be satisfied by putting an integer for $x$, then this integer, as the equation (1) shows, must necessarily be a divisor of $a_n$, and one has therefore only to substitute the various divisors of $a_n$, taken with positive and negative sign, into the equation by way of trial.

To the question of rational roots of an equation with integral coefficients there also comes back the general question whether an integral rational function
\[
f(x) = x^n + a_1 x^{n-1} + a_2 x^{n-2} + \cdots + a_{n-1} x + a_n,
\]
whose coefficients $a_1, a_2 \ldots a_n$ are integers, can be divisible by a rational function of lower degree $\nu$:
\[
\varphi(x) = x^\nu + \alpha_1 x^{\nu-1} + \alpha_2 x^{\nu-2} + \cdots + \alpha_{\nu-1} x + \alpha_\nu
\]
with rational coefficients $\alpha_1, \alpha_2 \ldots \alpha_\nu$. If $f(x)$ is divisible by $\varphi(x)$, then the quotient $\psi(x)$ is again an integral rational function with rational coefficients of degree $n - \nu$:
\[
\psi(x) = x^{n-\nu} + \beta_1 x^{n-\nu-1} + \cdots + \beta_{n-\nu-1} x + \beta_{n-\nu},
\]
and since $f(x)$ is also divisible by $\psi(x)$, we can, in
%CONT>>
